\section{Discussion and Limitations}\label{sec:discussion}

\paragraph{What would make this paper wrong.}
Each claim has a stated failure condition, and we would rather name them than defend them: A fails if \llm pools below $53.2\%$ or costs more per task than the API baseline; B fails if $196$ tokens are not worse than $576$/$1024$ on detail-dense views; C fails as a \emph{test} if the ladder is not re-measured, not if its span comes out small---a small span is a result about substitution.
% 中文：什么会证明本文错了。每条主张都有明示的失败条件：主张 A 失败于合并低于 53.2\% 或每任务成本高于 API 基线；主张 B 失败于 196 token 在细节密集视图上不差于 576/1024；主张 C 若未重测才算失败，跨度小不算——小跨度是关于「替代」的结果。

\paragraph{Limitations.}
First, the visual study covers PyAether views; SKILL and Tcl visuals are not measured, so any claim that vision helps ``EDA generally'' is out of scope here.
% 中文：局限。第一，视觉研究只覆盖 PyAether 视图；SKILL 与 Tcl 的视觉未测，因此「视觉对 EDA 普遍有用」的说法超出本文范围。
Second, $196$ tokens is the operating point of a \emph{general} encoder, chosen on loss and throughput rather than on EDA structure; a purpose-built encoder could move the threshold, and RQ5 measures the threshold rather than optimising it.
% 中文：第二，196 token 是「通用」编码器的工作点，依损失与吞吐选定而非依 EDA 结构；专用编码器可能挪动阈值，而 RQ5 是测量阈值而非优化阈值。
Third, the interaction between the RL reward and the anti-cheat hook is unverified: a policy optimised against assertions may learn to satisfy the hook rather than the task, and while we report a cheat rate beside every accuracy, we have not shown that the rate is stable under RL pressure.
% 中文：第三，RL 奖励与反作弊钩子的交互未验证：针对断言优化的策略可能学会满足钩子而非任务；虽然我们在每个精度旁报告作弊率，但未证明该速率在 RL 压力下稳定。
Fourth, the serving numbers are single-node; multi-tenant amortisation is modelled but not measured.
% 中文：第四，服务数字为单节点；多租户摊销只是建模，未实测。

\paragraph{If the 2B model collapses on SKILL.}
The subagent proposition would then hold only for PyAether---a toolchain with introspectable APIs, where execution can substitute for retrieval.
% 中文：若 2B 在 SKILL 上崩塌。届时 subagent 命题只对 PyAether 成立——那是一个 API 可自省、执行可替代检索的工具链。
That is a narrower but still publishable claim, and it would sharpen rather than weaken ZhuLong's toolchain-specific finding: component contributions differ by toolchain, so the \emph{deployable} subset of the harness differs by toolchain too.
% 中文：这是一个更窄但仍可发表的主张，且会强化（而非削弱）ZhuLong 的工具链特定结论：组件贡献因工具链而异，故框架的「可部署子集」也因工具链而异。
We pre-commit to reporting this boundary instead of averaging it away in the pooled number.
% 中文：我们预先承诺报告这一边界，而不是用合并数字把它抹平。
